\section{Theory and Hypotheses}

\subsection{From transaction price to public signal}

A prediction-market price is not self-interpreting. Institutions select a contract, choose a displayed price, attach a label, and place the number within a narrative. We describe this process as the manufacture of a public signal. Manufacture does not imply falsification; it denotes institutional selection, formatting, interpretation, and circulation.

We distinguish four levels of claim breadth. A transaction report states only that a contract traded at a particular price. A trader-belief claim interprets the price as the belief of market participants. An objective-probability claim removes the reference to participating traders and presents the number as the event probability. A public-representation claim further generalizes the number to public sentiment, consensus, or what a population believes. Each transition introduces assumptions about aggregation, probability interpretation, and representativeness.

\subsection{Institutional framing}

Prediction-market platforms have incentives to emphasize information production, trading, and participation. Media partners may also value a continuously updating numerical object that can be embedded in news coverage. Independent and regulatory sources face different incentives and professional routines. These institutional differences motivate the following hypotheses:

\begin{description}
\item[H1.] Partner-media and platform-owned texts exhibit higher epistemic-authority framing than independent-media and regulatory/critical texts.
\item[H2.] Partner-media and platform-owned texts are more likely to present prices as public sentiment, collective belief, or consensus.
\item[H3.] Partner-media and platform-owned texts exhibit less uncertainty disclosure and gambling/risk framing than regulatory/critical texts.
\end{description}

These are associational hypotheses. A partnership indicator does not randomly assign editorial framing.

\subsection{Disclosure and quality alignment}

A broad informational claim is better supported when audiences receive relevant context: who participates, how actively the contract trades, whether the spread is wide, whether the price is stale or volatile, and whether resolution rules are clear. Because such details complicate a simple data narrative, they may be omitted even when authority claims are strong.

\begin{description}
\item[H4.] Conditional on an objective-probability or public-representation claim, partner-media and platform-owned mentions are less likely to disclose market-quality or representativeness limitations.
\item[H5.] The positive relationship between epistemic-authority framing and contract-level market quality is weak or specification-sensitive.
\end{description}

H5 distinguishes a quality-selection account from a rhetorical-translation account. Under a strong quality-selection account, media use authoritative language primarily for deep, active, stable markets. Under a rhetorical-translation account, authoritative presentation is not reliably concentrated in the highest-quality contracts.

\subsection{Representational overreach}

We define representational overreach as a configuration in which the breadth and authority of a public claim exceed the evidentiary support made visible by market quality and disclosure. It is not synonymous with forecast error, deception, illegality, or gambling. A market can ultimately resolve correctly and still have been represented too broadly at the time of coverage.

The primary classification requires three conditions: (1) an objective-probability or public-representation claim; (2) low market quality under prespecified, platform-comparable criteria; and (3) no relevant limitation disclosure.

\begin{description}
\item[H6.] High-overreach configurations are disproportionately concentrated in partner-media and platform-owned mentions.
\end{description}

The analysis will also report a continuous score and a threshold specification curve. The transparent conjunction remains primary because it makes the normative and empirical ingredients auditable.
